\scanpage{36}
% archive foliation: 35
\uline{\unsure{Exponentielle|Exponentielles}} (en car. 0)

\margin{\deleted{\unsure{I)}}}
I) \uline{Généralités}

$A$ $\mathbb{Q}$-algèbre (ass. avec 1)

$N(A)$ ens des él. nilp. de $A$. On a donc $1+N(A)\subset A^*$ (ens des él. inv. de $A$)
\[
N(A)\xrightarrow{\exp}1+N(A)\qquad\exp x=\sum_{n\geq0}\frac{x^n}{n!}
\]
\uline{Th 1} 1) $\exp$ est bijective. L'application inverse est donnée par
\[
\log(1+y)=\sum_{n\geq1}(-1)^{n+1}\frac{y^n}{n}\qquad(y\in N(A))
\]
2) Si $x,y\in N(A)$ commutent, alors $x+y\in N(A)$ et
\[
\exp(x+y)=\exp x\ \exp y
\]
\margin{NB $\exp0=1$}
Si $u,v\in1+N(A)$ commutent, alors $uv\in1+N(A)$, et
\[
\log(uv)=\log u+\log v
\]
\margin{NB $\exp1=0$}
3) Soit $L\subset A$ un sous-$\mathbb{Q}$-espace vectoriel tel que $x\in L$, $n\in\mathbb{N}^*$ implique $x^n\in L$, \deleted{\unsure{Alors}} et $L$ formé d'éléments nilpotents. Alors
\[
x\longmapsto\exp x:\quad L\xrightarrow{\sim}1+L
\]
\uline{Remarque} Soit \deleted{\unsure{les conditions de 3), supposons que}} \inserted{$L\subset A$ un sous $\mathbb{Q}$-espace vectoriel de $A$} tel que $L\subset N(A)$, et tel que $x,y\in L\Rightarrow xy\in L$. Alors $L\xrightarrow{\sim}1+L$ par $3^\circ$, et de plus $1+L$ est un sous-groupe de $A^*$. D'où par transport de str une str. de groupe sur $L$. Elle est comm. \unsure{ssi|si} $L$ est commutative, i.e. si $x,y\in L\Rightarrow xy=yx$ --- \inserted{et dans ce cas cette structure est l'\uline{additive} de $L$.}
\margin{réciproque ? \unsure{par introduction|par l'introduction} \unsure{d'indéterminées} \unsure{voir} ---}

\uline{Exemple 1} Soient $K$ une $\mathbb{Q}$-alg. comm., $A$ une $K$-algèbre (ass. avec 1), $J$ un \deleted{\unsure{formé}} \inserted{nil} idéal de $K$. Alors $J.A$ satisfait les conditions de la Remarque, et on a donc
\[
JA\longrightarrow1+JA
\]
\uline{NB} Soient $K_0=K/J$, $A_0=A/JA=A\otimes_KK_0$, \inserted{$\unsure{\alpha_U|\alpha_A}:A\to A_0$ app. can.} Alors \deleted{\unsure{JA\longrightarrow}}
\[
JA=\{x\in A\mid\unsure{\alpha_U|\alpha_A}(x)=0\},\qquad1+JA=\{x\in A\mid\unsure{\alpha_U|\alpha_A}(x)=1\}
\]
\uline{Exemple 2} $K,J$ comme dans Exemple 1, \inserted{mais ici $J$ \uline{nilpotent},} $M$ un $K$-module,
\[
U=\operatorname{End}_K(M),\quad
\deleted{\unsure{U_0=U\otimes_KK_0\simeq U/JU}}\quad
\inserted{M_0=M\otimes_KK_0\simeq M/\unsure{JM}},
\]
$\alpha_M:M\longrightarrow M_0$ appl. can.; posons
\[
L=\{u\in U\unsure{=|:}\operatorname{End}_K(M)\mid u\otimes_KK_0=0\}
=\{u\in\unsure{M|U}\mid u(M)\subset JM\}
\]
alors
\[
1+L=\{u\in U=\operatorname{End}_K(M)\mid u\otimes_KK_0=\operatorname{id}_{M_0}\}.
\qquad\deleted{\unsure{\text{Ceci posé}}}
\]
Ceci posé, les conditions de la remarque sont satisfaites, donc on a
\[
\exp:L\xrightarrow{\sim}1+L
\]
